\section{Related Work}
\label{sec:related}

\subsection{Measurement Error in Causal Inference}

Classical measurement error theory shows that naive regression with mismeasured covariates yields attenuated coefficients~\cite{fuller1987measurement,carroll2006measurement}.
\citet{frazis2003estimating} derive exact bias formulas for ordinary least squares (OLS) with a misclassified binary regressor, showing attenuation proportional to the error rate, but restricted to average effects without subgroup heterogeneity.
In the econometrics literature, \citet{knox2020administrative} demonstrate that misclassification in administrative records biases estimates of racial disparities, and \citet{bound2001measurement} survey how binary misclassification attenuates regression coefficients.
Our Theorem~\ref{thm:bias_identity} generalizes the Frazis--Loewenstein scalar result to subgroup-heterogeneous error rates, revealing a qualitatively new product-form bias that can reverse effect signs across subgroups.
These works treat misclassification as a uniform nuisance; none models the interaction between heterogeneous error rates and heterogeneous treatment effects that arises when noisy labels serve as effect modifiers.

\subsection{Post-Prediction Causal Inference}

Several frameworks correct for machine-generated label noise in causal estimation.
PPI~\cite{angelopoulos2023ppi} and DSL~\cite{egami2023dsl} use a small expert-labeled sample to debias predictions, assuming homogeneous prediction quality across observations.
PPCI~\cite{sui2026ppci} extends this to conditional prediction correction, and ForestIV~\cite{bu2023forestiv} uses forest-based instrumental variables.
\citet{audinet2025benchmarking} benchmark these methods and report that debiasing can \emph{increase} MSE when prediction error is heterogeneous---a finding we explain: their population-level MSE increase is a direct consequence of the $\delta_s \cdot \Delta\tau$ product in Theorem~\ref{thm:bias_identity}, which shows that global correction amplifies bias whenever error-rate deviation interacts with effect heterogeneity.
All of these frameworks assume a single global error model; none addresses per-subgroup effect modifier misclassification, leaving a blind spot precisely where LLM annotation heterogeneity meets causal heterogeneity.

\subsection{Heterogeneous Treatment Effects}

The gap left by global error correction becomes consequential when corrected labels feed into HTE estimation.
Modern HTE estimators, including causal forests~\cite{wager2018causal}, double machine learning (DML)~\cite{chernozhukov2018dml}, and recursive partitioning~\cite{athey2016recursive}, assume error-free covariates.
When the effect modifier $Z$ is an LLM annotation, error rates vary substantially across demographic subgroups~\cite{ashwin2025bias}, violating this assumption.
\citet{imai2010causal} study treatment misclassification (errors in $T$, not $Z$), requiring instrumental variables; our setting preserves the treatment contrast but corrupts the heterogeneity structure, making confusion matrix inversion feasible.
Extended discussion appears in Appendix~\ref{app:related}.
